\section{Personas}

Al fine di calare nel mondo reale i tratti concettuali degli attori delineati in precedenza, sono state elaborate quattro \textit{Personas} che ne incarnano i background e gli scenari d'utilizzo tipici.

\subsection{Giustificazione del numero di Personas}

Si è scelto di adottare \textbf{quattro} Personas anziché tre (una per macro-attore) poiché l'analisi dei goal ha evidenziato che il macro-attore \textit{Turista Interno} racchiude in realtà due profili con bisogni antitetici: chi conosce a fondo il sito e cerca approfondimenti di nicchia, e chi invece, pur essendo del luogo, non lo ha mai visitato e necessita di un accompagnamento analogo a quello del turista esterno ma con tempistiche e motivazioni completamente differenti. Accorpare questi due profili in un'unica Persona avrebbe prodotto un personaggio internamente contraddittorio, incapace di guidare scelte progettuali univoche. D'altro canto, aggiungere ulteriori Personas (ad esempio distinguendo per fascia d'età o per competenza digitale) non avrebbe apportato differenze significative nei goal primari, bensì solo sfumature secondarie gestibili tramite personalizzazione dell'interfaccia. Quattro Personas rappresentano dunque il compromesso ottimale tra copertura degli scenari reali e maneggevolezza progettuale.

\subsection{Persona 1: Marco, il Turista Esterno}
\begin{itemize}
    \item \textbf{Età e Occupazione:} 46 anni, impiegato nel ramo assicurativo.
    \item \textbf{Background:} Proviene da una regione differente e si trova in viaggio turistico con la famiglia per il fine settimana; per lui è la prima volta in città.
    \item \textbf{Scenario d'uso:} Marco approda al complesso culturale senza avere familiarità con l'offerta interna. Dispone di appena un paio d'ore per districarsi nella visita dopodiché dovrà rincasare.
    \item \textbf{Obiettivi e Bisogni (Goal):} Avverte la necessità di un compagno di viaggio digitale per intercettare agilmente le \textit{highlight} (le gemme artistiche note) tramite descrizioni coincise, evitando sterili dispersioni logiche.
    \item \textbf{Frustrazioni:} Si inasprisce di fronte a densi blocchi di testo e si arrende qualora lo schema a video si riveli macchinoso negandogli lo sfruttamento repentino del ridotto capitale di tempo.
\end{itemize}

\subsection{Persona 2: Elena, la Turista Interna (Visitatrice Assidua)}
\begin{itemize}
    \item \textbf{Età e Occupazione:} 26 anni, laureanda presso la facoltà di Lettere e Beni Culturali.
    \item \textbf{Background:} È residente storicamente nel comprensorio cittadino. Lo ha esplorato materialmente in svariate edizioni e fiere.
    \item \textbf{Scenario d'uso:} Approfittando della calma di un mercoledì pomeriggio, varca le soglie per contemplare individualmente alcune rifiniture architettoniche recentemente restaurate.
    \item \textbf{Obiettivi e Bisogni (Goal):} Indaga su elementi solitamente evasi dalla placca ufficiale, quali storie occulte, retroscena artistici e percorsi fuori focolare. Non interroga per nulla le geolocalizzazioni perimetrali, essendo consapevole dello schema base, ma naviga con ardore traiettorie tematiche trasversali fornite dall'applicativo.
    \item \textbf{Frustrazioni:} Ripudia una banale ristesura enciclopedica Wikipedia e rimane profondamente svilita qualora constati un difetto di approfondimento, percependo il mezzo come sterile surrogato per "profani".
\end{itemize}

\subsection{Persona 3: Matteo, il Turista Interno (Prima Visita)}
\begin{itemize}
    \item \textbf{Età e Occupazione:} 34 anni, infermiere al policlinico.
    \item \textbf{Background:} Abitante periferico che orbita costantemente nei pressi del polo. Benché a conoscenza della reputazione strabiliante d'attrazione, rinvii d'inerzia lo hanno trattenuto dall'osservarlo nel corso della sua vita.
    \item \textbf{Scenario d'uso:} Stimolato da un incentivo feriale, in compagnia della logopedia di reparto, si risolve di affrontare la tanto rimandata full immersion.
    \item \textbf{Obiettivi e Bisogni (Goal):} È divorato dal duplice fabbisogno: basi accademiche macroscopiche combinato ad accessibilità emotiva di transito. Ricerca il "grande classico iconico", senza rinunciare ad accoppiamenti laterali godendosi una distesa andatura autoctona (lui il tempo lo ha dalla sua parte).
    \item \textbf{Frustrazioni:} Avverte lo scoramento quando incespica in appalti da accademia troppo cerebrali, o laddove si dia per assodato il fatto che gli inquilini del quartiere "debbano in automatico sapersi divincolare nei corridoi" mancando le opportune stesure logistiche.
\end{itemize}

\subsection{Persona 4: Giovanni, l'Impiegato del polo culturale}
\begin{itemize}
    \item \textbf{Età e Occupazione:} 55 anni, personale di salvaguardia e preposto ausiliario museale.
    \item \textbf{Background:} Detiene anzianità ventennale tra quelle gallerie; in virtù del suo mestiere assimila senza eguali perplessità e incagli turistici.
    \item \textbf{Scenario d'uso:} Nel vivo dell'operatività di turnazione, estrae lo smartphone impadronendosi di \textit{Culturia} al fine di sbloccare o coadiuvare la visita.
    \item \textbf{Obiettivi e Bisogni (Goal):} Si proietta in ottiche documentali o diagnostiche; ispeziona che gli accreditamenti a schermo mantengano convergenza col comparto fisico esposto. Adopera visivamente i tracciati colorati in app agitando il dispositivo per chiarificare la sosta d'emergenza o i ristori idrici ai forestieri.
    \item \textbf{Frustrazioni:} Rigetta pesantemente l'uso allorché si avveda d'informazioni inautentiche diramate al cittadino, così come paleserà forte disagio nel combattere con menù inestricabili nell'istante perentorio di somministrare aiuti sul posto, di fronte alla clientela esigente.
\end{itemize}
